\section{Introduction}

Deep Reinforcement Learning (RL) has become the dominant paradigm for developing competent agents in challenging environments. Most notably, deep RL algorithms have achieved impressive performance in a multitude of arcade \citep{dqn, muzero, dreamerv2}, real-time strategy \citep{alphastar, dota2}, board \citep{alphago, alphazero, muzero} and imperfect information \citep{pog, rebel} games. However, a common drawback of these methods is their extremely low sample efficiency. Indeed, experience requirements range from months of gameplay for DreamerV2 \citep{dreamerv2} in Atari 2600 games \citep{bellemare13arcade} to thousands of years for OpenAI Five in Dota2 \citep{dota2}. While some environments can be sped up for training agents, real-world applications often cannot. Besides, additional cost or safety considerations related to the number of environmental interactions may arise \citep{aisecurity}. Hence, sample efficiency is a necessary condition to bridge the gap between research and the deployment of deep RL agents in the wild.

Model-based methods \citep{sutton} constitute a promising direction towards data efficiency. Recently, world models were leveraged in several ways: pure representation learning \citep{spr}, lookahead search \citep{muzero, efficientzero}, and learning in imagination \citep{worldmodels, simple, dreamerv1, dreamerv2}. The latter approach is particularly appealing because training an agent inside a world model frees it from sample efficiency constraints. Nevertheless, this framework relies heavily on accurate world models since the policy is purely trained in imagination. In a pioneering work, \citet{worldmodels} successfully built imagination-based agents in toy environments. SimPLe recently showed promise in the more challenging Atari 100k benchmark \citep{simple}. Currently, the best Atari agent learning in imagination is DreamerV2 \citep{dreamerv2}, although it was developed and evaluated with two hundred million frames available, far from the sample-efficient regime. Therefore, designing new world model architectures, capable of handling visually complex and partially observable environments with few samples, is key to realize their potential as surrogate training grounds.

The Transformer architecture \citep{vaswani} is now ubiquitous in Natural Language Processing \citep{bert, radford_gpt2, brown_gpt3, t5}, and is also gaining traction in Computer Vision \citep{vit, masked_autoencoders_he}, as well as in Offline Reinforcement Learning \citep{trajectory_transformer, decision_transformer}. In particular, the GPT \citep{gpt1, radford_gpt2, brown_gpt3} family of models delivered impressive results in language understanding tasks. Similarly to world models, these attention-based models are trained with high-dimensional signals and a self-supervised learning objective, thus constituting ideal candidates to simulate an environment.

Transformers particularly shine when they operate over sequences of discrete tokens \citep{bert, brown_gpt3}. For textual data, there are simple ways \citep{wordpiece, sentencepiece} to build a vocabulary, but this conversion is not straightforward with images. A naive approach would consist in treating pixels as image tokens, but standard Transformer architectures scale quadratically with sequence length, making this idea computationally intractable. To address this issue, \textsc{vqgan} \citep{vqgan} and \textsc{dall-e} \citep{ramesh} employ a discrete autoencoder \citep{vqvae} as a mapping from raw pixels to a much smaller amount of image tokens. Combined with an autoregressive Transformer, these methods demonstrate strong unconditional and conditional image generation capabilities. Such results suggest a new approach to design world models. 

\definecolor{autoencodercolor}{rgb}{0.067,0.466,0.2}
\definecolor{gptcolor}{rgb}{0.2, 0.133, 0.533}
\definecolor{policycolor}{rgb}{0.667,0.266,0.6}

\begin{figure}[t]
\center
\begin{tikzpicture}[font=\scriptsize, scale=1.15]

  \node (z0) at (0, 0) {$z^1_0, \dots, z^K_0$};
  \node (a0) [right=3pt of z0] {$a_0$};
  \node (x0) [above=1.25cm of z0] {$x_0$};
  \node (xh0) [below=1.25cm of z0] {$\hat{x}_0$};

  \node[draw, inner sep=0pt, very thin, densely dashed, fit=(z0)] (bz0) {};
  \node[draw, inner sep=4pt, very thin, fit=(bz0) (a0)] (az0) {};

  \node (z1) at (3.65, 0) {$\hat{z}^1_1, \dots, \hat{z}^K_1$};
  \node (a1) [right=3pt of z1] {$a_1$};
  \node (xh1) [below=1.25cm of z1] {$\hat{x}_1$};
  \node (dh1) [above=1.0cm of z1.west,xshift=8pt] {$\hat{d}_0$};
  \node (rh1) [above=0pt of dh1] {$\hat{r}_0$};
  \node[draw, inner sep=0pt, very thin, densely dashed, fit=(z1)] (bz1) {};
  \node[draw, inner sep=4pt, very thin, fit=(az0) (bz1) (a1)] (az1) {};

  \node[minimum width=0.5cm] (dots) at (6.75, 0) {$\ \dots$};
  \node[draw, inner sep=4pt, very thin, fit=(az1) (dots)] (azd) {};
  \node (dh2) [above=1.0cm of dots.west,xshift=8pt] {${\hat{d}_1}$};
  \node (rh2) [above=0pt of dh2] {${\hat{r}_1}$};

  \node (zt) at (9.2, 0) {$\hat{z}^1_s, \dots, \hat{z}^K_s$};
  \node (at) [right=3pt of zt] {$a_s$};
  \node (xht) [below=1.25cm of zt] {$\hat{x}_s$};
  \node (dht) [above=1.0cm of zt.west,xshift=8pt] {$\hat{d}_{s-1}$};
  \node (rht) [above=0pt of dht] {$\hat{r}_{s-1}$};
  \node[draw, inner sep=0pt, very thin, densely dashed, fit=(zt)] (bzt) {};

  \coordinate (az01) at ($(az0.east)!0.6!(bz1.west)$);
  \draw[very thick,gptcolor] (az0) -- (az01) node[midway,above] {$G$};
  \draw[->,very thick,gptcolor] (az01) -- (bz1);
  \draw[->,very thick,gptcolor] (az01) |- (rh1);
  \draw[->,very thick,gptcolor] (az01 |- dh1) -- (dh1);

  \coordinate (az02) at ($(az1.east)!0.6!(dots.west)$);
  \draw[very thick,gptcolor] (az1) -- (az02) node[midway,above] {$G$};
  \draw[->,very thick,gptcolor] (az02) -- (dots);
  \draw[->,very thick,gptcolor] (az02) |- (rh2);
  \draw[->,very thick,gptcolor] (az02 |- dh2) -- (dh2);

  \coordinate (az0t) at ($(azd.east)!0.6!(bzt.west)$);
  \draw[very thick,gptcolor] (azd) -- (az0t) node[midway,above] {$G$};
  \draw[->,very thick,gptcolor] (az0t) -- (bzt);
  \draw[->,very thick,gptcolor] (az0t) |- (rht);
  \draw[->,very thick,gptcolor] (az0t |- dht) -- (dht);

  %% \draw[->,thick] (azd) -- (bzt) node[midway,above] {$G$};

  \draw[->,very thick,autoencodercolor] (x0) -- (bz0) node[midway,right,yshift=5pt] {$E$};
  \draw[->,very thick,autoencodercolor] (bz0) -- (xh0) node[midway,right,yshift=-5pt] {$D$};
  \draw[->,very thick,policycolor] (xh0) -- (xh0 -| a0) node[midway,below] {$\pi$} -- (a0);
  \draw[->,very thick,autoencodercolor] (bz1) -- (xh1) node[midway,right,yshift=-5pt] {$D$};
  \draw[->,very thick,policycolor] (xh1) -- (xh1 -| a1) node[midway,below] {$\pi$} -- (a1);
  \draw[->,very thick,autoencodercolor] (bzt) -- (xht) node[midway,right,yshift=0pt] {$D$};
  \draw[->,very thick,policycolor] (xht) -- (xht -| at) node[midway,below] {$\pi$} -- (at);

  \draw[thick,decoration={brace,mirror,amplitude=0.25cm},decorate] ($(z0)-(0.8, 2)$) -- ++(2.3, 0) node[midway,below,yshift=-8pt] {$t=0$};
  \draw[thick,decoration={brace,mirror,amplitude=0.25cm},decorate] ($(z1)-(0.8, 2)$) -- ++(2.3, 0) node[midway,below,yshift=-8pt] {$t=1$};
  \node[minimum width=1.5cm] (dots2) at ($(z1)!0.5!(zt)+(0.25, -2.1)$) {$\dots$};
  \draw[thick,decoration={brace,mirror,amplitude=0.25cm},decorate] ($(zt)-(0.8, 2)$) -- ++(2.3, 0) node[midway,below,yshift=-8pt] {$t=s$};

\end{tikzpicture}

\caption[Unrolling imagination over time.]{Unrolling imagination over time. This figure shows the policy $\pi$, depicted with purple arrows, taking a sequence of actions in imagination. The green arrows correspond to the encoder $E$ and the decoder $D$ of a discrete autoencoder, whose task is to represent frames in its learnt symbolic language. The backbone $G$ of the world model is a GPT-like Transformer, illustrated with blue arrows. For each action that the policy $\pi$ takes, $G$ simulates the environment dynamics, by autoregressively unfolding new frame tokens that $D$ can decode. $G$ also predicts a reward and a potential episode termination. More specifically, an initial frame $x_0$ is encoded with $E$ into tokens $z_0 = (z_0^1, \dots, z_0^K) = E(x_0)$. The decoder $D$ reconstructs an image $\hat{x}_0 = D(z_0)$, from which the policy $\pi$ predicts the action $a_0$. From $z_0$ and $a_0$, $G$ predicts the reward $\hat{r}_0$, episode termination $\hat{d}_0 \in \{0, 1\}$, and in an autoregressive manner $\hat{z}_1 = (\hat{z}_1^1, \dots, \hat{z}_1^K)$, the tokens for the next frame. A dashed box indicates image tokens for a given time step, whereas a solid box represents the input sequence of $G$, \textit{i.e.} $(z_0, a_0)$ at $t = 0$,  $(z_0, a_0, \hat{z}_1, a_1)$ at $t = 1$, etc. The policy $\pi$ is purely trained with imagined trajectories, and is only deployed in the real environment to improve the world model $(E, D, G)$.} 

\vspace{-0.025\linewidth}
\label{fig:method}
\end{figure}

In the present work, we introduce \textsc{iris} (Imagination with auto-Regression over an Inner Speech), an agent trained in the imagination of a world model composed of a discrete autoencoder and an autoregressive Transformer. \textsc{iris} learns behaviors by accurately simulating millions of trajectories. Our approach casts dynamics learning as a sequence modeling problem, where an autoencoder builds a language of image tokens and a Transformer composes that language over time. With minimal tuning, \textsc{iris} outperforms a line of recent methods \citep{simple, rainbow, curl, drq, spr} for sample-efficient RL in the Atari 100k benchmark \citep{simple}. After only two hours of real-time experience, it achieves a mean human normalized score of 1.046, and reaches superhuman performance on 10 out of 26 games. We describe \textsc{iris} in Section \ref{sec:method} and present our results in Section \nolinebreak \ref{section:experiments}.